\documentclass[12pt, a4paper]{article}

% Paquetes requeridos
\usepackage[utf8]{inputenc}
\usepackage[T1]{fontenc}
\usepackage[spanish]{babel}
\usepackage{geometry}
\usepackage{amsmath}
\usepackage[american]{circuitikz}
\usetikzlibrary{babel} % Evita conflictos entre las flechas de circuitikz y babel-spanish

% Configuracion de la pagina
\geometry{a4paper, margin=2.5cm}

\begin{document}

\section*{Solución del Ejemplo 4.11: Análisis Nodal}

\subsection*{Planteamiento del Problema}
Considere el circuito eléctrico de la Figura \ref{fig:sol-ejm-4-11}. Determine las expresiones algebraicas para las tensiones en los nodos $a$, $c$ y $d$ aplicando análisis nodal, considerando los siguientes componentes:
\begin{itemize}
    \item Una fuente de corriente independiente $I$ y una resistencia $R_1$ conectadas en paralelo entre los nodos $a$ y $b$.
    \item Una resistencia $R_2$ entre los nodos $a$ y $c$, en la cual se define una caída de tensión $V_x$.
    \item Una resistencia $R_3$ conectada entre los nodos $c$ y $d$.
    \item Una fuente dependiente de tensión conectada entre los nodos $a$ y $d$, cuyo valor es $2V_x$.
    \item Una fuente dependiente de corriente conectada entre los nodos $c$ y $b$, cuyo valor es $2i_x$.
    \item Una fuente independiente de tensión de valor $V$ entre los nodos $d$ y $b$. La variable de control $i_x$ se define como la corriente que fluye desde el nodo $b$ hacia el nodo $d$ a través de esta fuente.
\end{itemize}

\begin{figure}[h!]
    \centering
    \begin{circuitikz}[american, scale=1.2, transform shape]
        
        % Línea inferior (Nodo común 'b')
        \draw (-2,0) -- (6,0);
        
        % Definición de Nodos principales a, c, d
        \node[circ] (A) at (0,3) {};
        \node[circ] (C) at (3,3) {};
        \node[circ] (D) at (6,3) {};
        
        % Etiquetas de los nodos
        \node[above left]  at (A) {a};
        \node[above]       at (C) {c};
        \node[above right] at (D) {d};
        \node[below]       at (3,0) {b (Ref)};
        
        % Puntos de unión en el riel inferior 'b'
        \node[circ] at (-2,0) {};
        \node[circ] at (0,0) {};
        \node[circ] at (3,0) {};
        \node[circ] at (6,0) {};

        % Componentes
        \draw (-2,0) to[I, l=$I$] (-2,3) -- (0,3);
        \draw (0,3) to[R, l=$R_1$] (0,0);
        \draw (0,3) to[R, l=$R_2$, v=$V_x$] (3,3);
        \draw (3,3) to[R, l=$R_3$] (6,3);
        \draw (0,3) -- (0,5) to[cV, l=$2V_x$] (6,5) -- (6,3);
        \draw (3,3) to[cI, l=$2i_x$] (3,0);
        \draw (6,3) to[V, l_=$V$, i<=$i_x$] (6,0);

    \end{circuitikz}
    \caption{Circuito del Ejemplo 4-11.}
    \label{fig:sol-ejm-4-11}
\end{figure}

\subsection*{Desarrollo Matemático}

\textbf{Paso 1: Definir el nodo de referencia e identificar tensiones fijadas.} \\
Asignamos el nodo $b$ como referencia ($V_b = 0\,\text{V}$). Dado que existe una fuente de tensión independiente $V$ conectada entre el nodo $d$ y el nodo $b$, la tensión en el nodo $d$ es directamente:
\begin{equation}
    V_d = V
\end{equation}

\textbf{Paso 2: Ecuación de restricción de la fuente dependiente de tensión.} \\
La fuente dependiente de tensión está conectada entre $a$ y $d$, imponiendo que:
\begin{equation}
    V_a - V_d = 2V_x
\end{equation}
Sabemos que $V_x$ es la caída de tensión en $R_2$, es decir, $V_x = V_a - V_c$. Sustituyendo esto y $V_d = V$:
\begin{align*}
    V_a - V &= 2(V_a - V_c) \\
    V_a - V &= 2V_a - 2V_c \\
    V_a &= 2V_c - V \tag{3} \label{eq:Va}
\end{align*}

\textbf{Paso 3: Hallar una expresión para $i_x$.} \\
La variable $i_x$ es la corriente que sube por la fuente $V$. Para evitar cálculos complejos con la corriente de la fuente $2V_x$, aplicamos la Ley de Corrientes de Kirchhoff (LKC) en el nodo de referencia $b$ (sumatoria de corrientes que salen = sumatoria de corrientes que entran):
\begin{align*}
    \text{Salen de } b: \quad & I + i_x \\
    \text{Entran a } b: \quad & \frac{V_a}{R_1} + 2i_x
\end{align*}
Igualando y despejando $i_x$:
\begin{equation}
    I + i_x = \frac{V_a}{R_1} + 2i_x \quad \Rightarrow \quad i_x = I - \frac{V_a}{R_1} \tag{4} \label{eq:ix}
\end{equation}

\textbf{Paso 4: LKC en el nodo $c$.} \\
Aplicamos LKC en el nodo $c$ (sumatoria de corrientes que salen = 0):
\begin{equation}
    \frac{V_c - V_a}{R_2} + \frac{V_c - V_d}{R_3} + 2i_x = 0
\end{equation}
Sustituimos $V_d = V$, la Ecuación \eqref{eq:Va} ($V_a = 2V_c - V$) y la Ecuación \eqref{eq:ix} ($i_x = I - \frac{2V_c - V}{R_1}$):
\begin{equation*}
    \frac{V_c - (2V_c - V)}{R_2} + \frac{V_c - V}{R_3} + 2\left(I - \frac{2V_c - V}{R_1}\right) = 0
\end{equation*}
\begin{equation*}
    \frac{V - V_c}{R_2} + \frac{V_c - V}{R_3} + 2I - \frac{4V_c - 2V}{R_1} = 0
\end{equation*}
Agrupando los términos que contienen $V_c$ a un lado y los términos independientes ($V$, $I$) al otro, obtenemos la tensión en el nodo $c$:
\begin{equation}
    V_c = \frac{V \left( \frac{2}{R_1} + \frac{1}{R_2} - \frac{1}{R_3} \right) + 2I}{\frac{4}{R_1} + \frac{1}{R_2} - \frac{1}{R_3}}
\end{equation}
Con este resultado de $V_c$, la tensión $V_a$ se calcula fácilmente usando la Ecuación \eqref{eq:Va}.
\end{document}